% ============================================================
% CHAPITRE 2 — ÉTAT DE L'ART
% ============================================================

\chapter{État de l'Art}

\section*{Plan}

\begin{tabular}{r l r}
\textbf{1} & \textbf{Les modèles de langage de grande taille (LLM)} & \textbf{\pageref{sec:llm}} \\
\textbf{2} & \textbf{Gouvernance et observabilité des LLM} & \textbf{\pageref{sec:gouvernance}} \\
\textbf{3} & \textbf{Architecture API Gateway et pattern Reverse Proxy} & \textbf{\pageref{sec:apigateway}} \\
\textbf{4} & \textbf{Séries temporelles et bases de données temporelles} & \textbf{\pageref{sec:timeseries}} \\
\textbf{5} & \textbf{Cache sémantique et embeddings vectoriels} & \textbf{\pageref{sec:embeddings}} \\
\textbf{6} & \textbf{Architectures multi-tenant SaaS} & \textbf{\pageref{sec:multitenant}} \\
\textbf{7} & \textbf{Sécurité des API~: JWT, hachage et bonnes pratiques} & \textbf{\pageref{sec:securite}} \\
\textbf{8} & \textbf{Détection d'anomalies par apprentissage automatique} & \textbf{\pageref{sec:anomaly}} \\
\textbf{9} & \textbf{Prévision de séries temporelles} & \textbf{\pageref{sec:forecasting}} \\
\textbf{10} & \textbf{Architecture événementielle et tâches asynchrones} & \textbf{\pageref{sec:eventdriven}} \\
\textbf{11} & \textbf{Conteneurisation et intégration continue} & \textbf{\pageref{sec:conteneurisation}} \\
\textbf{12} & \textbf{Solutions existantes et analyse comparative} & \textbf{\pageref{sec:comparatif}} \\
\end{tabular}

\section*{Introduction}

Ce chapitre dresse un panorama exhaustif des fondements théoriques et technologiques qui sous-tendent notre plateforme de gouvernance des LLM. Nous explorons successivement les modèles de langage de grande taille et l'architecture Transformer qui les rend possibles, les enjeux de gouvernance et d'observabilité en production, l'architecture des API Gateways et des reverse proxies, les systèmes de gestion de séries temporelles, les embeddings vectoriels et le cache sémantique, les architectures multi-tenant SaaS, la sécurité des API par JWT et hachage cryptographique, les algorithmes de détection d'anomalies par apprentissage automatique, les méthodes de prévision de séries temporelles, les paradigmes événementiels et les tâches asynchrones, la conteneurisation et l'intégration continue, et enfin les solutions existantes sur le marché. Cette revue de la littérature, s'appuyant sur des sources scientifiques et des standards industriels, permet de situer notre contribution par rapport à l'état actuel des connaissances et de justifier rigoureusement chacun des choix techniques opérés dans les sprints de développement.

% ────────────────────────────────────────────────────────────
\section{Les modèles de langage de grande taille (LLM)}
\label{sec:llm}
% ────────────────────────────────────────────────────────────

\subsection{Définition et principes fondamentaux}

Un modèle de langage de grande taille (\textit{Large Language Model}, LLM) est un réseau de neurones profond pré-entraîné sur des corpus textuels massifs, capable de comprendre et de générer du langage naturel avec une fluidité remarquable. Ces modèles reposent sur l'apprentissage de distributions de probabilité sur des séquences de tokens --- unités élémentaires de texte correspondant à des mots, des sous-mots ou des caractères selon le schéma de tokenisation employé~\cite{brown2020}.

La puissance des LLM réside dans leur capacité à capturer des relations sémantiques complexes et à effectuer du \textbf{raisonnement en contexte} (\textit{in-context learning}), c'est-à-dire à s'adapter à de nouvelles tâches sans réentraînement, uniquement à partir d'exemples fournis dans le prompt~\cite{brown2020}. Cette propriété émergente, observée à partir d'une certaine échelle de paramètres, a transformé le paradigme de l'intelligence artificielle appliquée.

\subsection{L'architecture Transformer}

L'architecture \textbf{Transformer}, introduite par Vaswani \textit{et al.} en 2017~\cite{vaswani2017}, constitue le socle technique de tous les LLM modernes. Contrairement aux architectures récurrentes (RNN, LSTM) qui traitent les séquences de manière séquentielle, le Transformer s'appuie sur un mécanisme d'\textbf{attention multi-tête} (\textit{multi-head self-attention}) qui permet de modéliser les dépendances à longue distance de manière parallélisée.

Le mécanisme d'attention peut être formalisé comme suit~\cite{vaswani2017}~:

\begin{equation}
\text{Attention}(Q, K, V) = \text{softmax}\left(\frac{QK^T}{\sqrt{d_k}}\right)V
\end{equation}

\noindent où $Q$ (Query), $K$ (Key) et $V$ (Value) sont des projections linéaires de l'entrée, et $d_k$ est la dimension des clés, servant de facteur de normalisation. Cette formulation permet à chaque position de la séquence d'accéder directement à toutes les autres positions, éliminant ainsi le goulot d'étranglement séquentiel inhérent aux architectures récurrentes.

L'architecture Transformer se compose de deux blocs principaux~:

\begin{itemize}[noitemsep]
    \item l'\textbf{encodeur}~: transforme la séquence d'entrée en une représentation contextuelle dense~;
    \item le \textbf{décodeur}~: génère la séquence de sortie token par token, en s'appuyant sur la représentation produite par l'encodeur et sur les tokens déjà générés.
\end{itemize}

Les LLM modernes adoptent généralement une architecture de type \textbf{décodeur seul} (\textit{decoder-only}), où le modèle est entraîné de manière auto-régressive à prédire le token suivant dans une séquence~\cite{brown2020}.

\begin{figure}[H]
    \centering
    \fbox{\parbox{10cm}{\centering\vspace{3cm}\textit{Schéma de l'architecture Transformer}\vspace{3cm}}}
    \caption{Architecture Transformer originale~\cite{vaswani2017}}
    \label{fig:transformer}
\end{figure}

\subsection{Évolution historique des LLM}

L'évolution des modèles de langage peut être retracée à travers plusieurs jalons fondamentaux~:

\begin{itemize}[noitemsep]
    \item \textbf{GPT-1} (2018, OpenAI)~: premier modèle démontrant l'efficacité du pré-entraînement non supervisé suivi d'un fine-tuning supervisé, avec 117 millions de paramètres~;
    \item \textbf{BERT} (2018, Google)~: modèle bidirectionnel exploitant l'encodeur du Transformer, révolutionnant la compréhension du langage naturel~;
    \item \textbf{GPT-3} (2020, OpenAI)~: avec 175 milliards de paramètres, ce modèle a démontré des capacités émergentes de \textit{few-shot learning}~\cite{brown2020}~;
    \item \textbf{ChatGPT} (2022, OpenAI)~: fine-tuné par apprentissage par renforcement à partir de retours humains (RLHF), il a démocratisé l'usage des LLM~;
    \item \textbf{GPT-4} (2023, OpenAI)~: modèle multimodal aux performances proches de l'expertise humaine sur de nombreux benchmarks~\cite{openai2023}~;
    \item \textbf{LLaMA} (2023, Meta)~: famille de modèles open-source performants, contribuant à la démocratisation de la recherche~\cite{touvron2023}~;
    \item \textbf{Claude} (2023--2024, Anthropic)~: modèle axé sur la sécurité et l'alignement, avec des fenêtres de contexte étendues~\cite{anthropic2024}.
\end{itemize}

\subsection{Tokenisation et facturation}

La \textbf{tokenisation} est le processus de décomposition du texte en unités élémentaires appelées \textbf{tokens}. La plupart des LLM modernes utilisent des algorithmes de tokenisation sous-mot tels que \textbf{BPE} (\textit{Byte Pair Encoding}) ou \textbf{SentencePiece}. La facturation des services LLM s'effectue au token, avec des tarifs différenciés entre les tokens d'entrée (\textit{input tokens}) et les tokens de sortie (\textit{output tokens}). Cette structure tarifaire rend la maîtrise de la consommation de tokens essentielle pour les entreprises, justifiant l'existence d'une plateforme de monitoring dédiée.


% ────────────────────────────────────────────────────────────
\section{Gouvernance et observabilité des LLM}
\label{sec:gouvernance}
% ────────────────────────────────────────────────────────────

\subsection{Les défis de la production LLM}

Le déploiement de modèles de langage en production soulève des défis spécifiques que Bommasani \textit{et al.}~\cite{bommasani2022} et Chen \textit{et al.}~\cite{chen2024llmops} regroupent sous le terme de \textbf{LLMOps} (\textit{Large Language Model Operations}). Ces défis couvrent~:

\begin{itemize}[noitemsep]
    \item la \textbf{maîtrise des coûts}~: la facturation au token peut engendrer des coûts exponentiels si la consommation n'est pas surveillée et plafonnée~;
    \item la \textbf{traçabilité}~: chaque appel LLM doit être journalisé avec ses métadonnées (modèle utilisé, tokens consommés, latence, statut) pour l'audit et l'optimisation~;
    \item la \textbf{gouvernance}~: des règles de contrôle d'accès, de quotas et de budgets doivent être appliquées en temps réel~;
    \item la \textbf{sécurité}~: la protection des clés API, la prévention des injections de prompts et l'isolation des données entre tenants sont des impératifs~;
    \item l'\textbf{observabilité}~: la mise à disposition de métriques, de tableaux de bord et d'alertes est nécessaire pour une exploitation sereine.
\end{itemize}

\subsection{Le concept de proxy LLM}

Un \textbf{proxy LLM} est un composant architectural positionné entre les applications clientes et les fournisseurs de modèles. Il agit comme un intermédiaire transparent qui intercepte les requêtes, applique des politiques de gouvernance, journalise les interactions et route les appels vers le fournisseur approprié. Ce pattern architectural, inspiré des API Gateways traditionnelles~\cite{fielding2000}, constitue le c\oe ur de notre plateforme.

\begin{figure}[H]
    \centering
    \fbox{\parbox{12cm}{\centering\vspace{3cm}\textit{Architecture d'un proxy LLM avec gouvernance}\vspace{3cm}}}
    \caption{Principe du proxy LLM intelligent}
    \label{fig:llm_proxy}
\end{figure}


% ────────────────────────────────────────────────────────────
\section{Architecture API Gateway et pattern Reverse Proxy}
\label{sec:apigateway}
% ────────────────────────────────────────────────────────────

\subsection{Définition et rôle architectural}

Dans le contexte des architectures distribuées modernes, l'\textbf{API Gateway} désigne un composant qui sert de point d'entrée unique pour l'ensemble des requêtes adressées à un système composé de multiples services~\cite{richardson2018}. Richardson définit ce pattern comme un intermédiaire qui découple les clients de la topologie interne du système, en assurant des fonctions transversales telles que le routage, la composition de réponses, l'authentification et la limitation de débit.

Ce concept s'inscrit dans la lignée du pattern \textbf{Reverse Proxy}, formalisé dès les premières spécifications HTTP~\cite{fielding2000}. Contrairement à un proxy classique qui agit au nom du client, un reverse proxy se positionne côté serveur et reçoit les requêtes entrantes pour les redistribuer vers les services appropriés, tout en masquant l'infrastructure sous-jacente. Nginx, utilisé dans notre plateforme, est l'implémentation de référence de ce pattern, offrant des fonctionnalités natives de terminaison TLS, de compression et de \textit{rate limiting}~\cite{nginx2024}.

\subsection{Patterns de résilience}

Les systèmes distribués sont intrinsèquement exposés aux défaillances partielles. Pour y répondre, la littérature en ingénierie logicielle a formalisé plusieurs patterns de résilience que Richardson~\cite{richardson2018} et Nygard décrivent dans le contexte des microservices~:

\begin{itemize}[noitemsep]
    \item le \textbf{Circuit Breaker}~: inspiré de l'électrotechnique, ce pattern surveille le taux d'erreur d'un service aval et <<~ouvre le circuit~>> temporairement lorsqu'un seuil est dépassé, évitant ainsi la propagation en cascade des défaillances~;
    \item le \textbf{Rate Limiting}~: mécanisme de limitation du nombre de requêtes par unité de temps, typiquement implémenté par un compteur à fenêtre glissante dans un stockage en mémoire tel que Redis~;
    \item le \textbf{Retry with Backoff}~: stratégie de réessai avec délai exponentiel croissant, permettant de gérer les erreurs transitoires sans surcharger le service défaillant~;
    \item le \textbf{Bulkhead}~: isolation des ressources par domaine fonctionnel, empêchant qu'un composant défaillant ne monopolise l'ensemble des ressources du système.
\end{itemize}

\subsection{Application à notre plateforme}

Notre plateforme implémente le pattern API Gateway à deux niveaux. Le premier niveau est assuré par \textbf{Nginx}, qui agit comme reverse proxy en entrée de l'infrastructure, gérant la terminaison TLS, le \textit{rate limiting} global et le routage vers les services internes. Le second niveau est le \textbf{Gateway LLM}, un composant applicatif développé en FastAPI, qui intercepte chaque appel destiné aux fournisseurs LLM pour y appliquer la gouvernance, la journalisation et le cache sémantique avant de router la requête vers le provider sélectionné.


% ────────────────────────────────────────────────────────────
\section{Séries temporelles et bases de données temporelles}
\label{sec:timeseries}
% ────────────────────────────────────────────────────────────

\subsection{Définition et propriétés}

Une \textbf{série temporelle} est une séquence ordonnée de mesures indexées par le temps, formellement représentée par $\{(t_i, v_i)\}_{i=1}^{n}$ où $t_i$ désigne l'instant d'observation et $v_i$ la valeur mesurée~\cite{jensen2017}. Les séries temporelles possèdent des propriétés caractéristiques qui les distinguent des données relationnelles classiques~:

\begin{itemize}[noitemsep]
    \item l'\textbf{append-only}~: les nouvelles données sont exclusivement insérées en fin de série, les mises à jour étant exceptionnelles~;
    \item la \textbf{cardinalité temporelle}~: le volume croît linéairement avec le temps, pouvant atteindre plusieurs milliards d'enregistrements~;
    \item la \textbf{dépendance temporelle}~: les observations consécutives sont généralement corrélées, propriété exploitée par les algorithmes de détection d'anomalies~;
    \item la \textbf{granularité variable}~: les données sont souvent agrégées à différentes résolutions (minute, heure, jour, mois) pour des besoins analytiques distincts.
\end{itemize}

\subsection{Systèmes de gestion de séries temporelles}

Jensen, Pedersen et Thomsen~\cite{jensen2017} proposent une taxonomie des systèmes de gestion de séries temporelles (\textit{Time Series Management Systems}, TSMS) en distinguant trois grandes familles~:

\begin{table}[H]
\centering
\caption{Comparaison des approches de stockage de séries temporelles}
\label{tab:tsms_comparaison}
\begin{tabularx}{\textwidth}{|l|X|X|l|}
\hline
\textbf{Approche} & \textbf{Principe} & \textbf{Exemple} & \textbf{Cas d'usage} \\
\hline
SGBDR étendu & Extension d'un SGBDR classique avec des types et opérateurs temporels & TimescaleDB & Analytics, facturation \\
\hline
Base native & Base conçue spécifiquement pour les séries temporelles & InfluxDB & Monitoring infra \\
\hline
Stockage clé-valeur & Paires (timestamp, valeur) dans un stockage distribué & Prometheus & Métriques système \\
\hline
\end{tabularx}
\end{table}

\subsection{TimescaleDB et le concept de hypertable}

\textbf{TimescaleDB} est une extension open-source de PostgreSQL qui transforme une table relationnelle classique en une \textbf{hypertable}~\cite{timescaledb2024}. Une hypertable est une abstraction qui partitionne automatiquement les données en \textit{chunks} selon la dimension temporelle, offrant des performances d'insertion et de requête adaptées aux séries temporelles tout en préservant la compatibilité SQL complète de PostgreSQL.

Les avantages principaux de cette approche sont~:

\begin{itemize}[noitemsep]
    \item la \textbf{compression automatique}~: les chunks anciens sont compressés avec un ratio pouvant atteindre 95\,\%, réduisant considérablement l'empreinte disque~;
    \item les \textbf{agrégations continues} (\textit{continuous aggregates})~: des vues matérialisées qui se mettent à jour automatiquement à chaque nouvel enregistrement~;
    \item le \textbf{partitionnement transparent}~: les requêtes SQL standard bénéficient du partitionnement sans modification du code applicatif~;
    \item la \textbf{rétention de données}~: des politiques de suppression automatique des données anciennes, essentielles pour la conformité réglementaire.
\end{itemize}

Dans notre plateforme, la table \texttt{llm\_token\_log} --- qui enregistre chaque appel LLM avec ses métadonnées --- est configurée en hypertable partitionnée par la colonne \texttt{timestamp}, permettant de gérer efficacement un flux continu d'événements à haute fréquence tout en maintenant des performances de requête compatibles avec les tableaux de bord temps réel.


% ────────────────────────────────────────────────────────────
\section{Cache sémantique et embeddings vectoriels}
\label{sec:embeddings}
% ────────────────────────────────────────────────────────────

\subsection{Représentations vectorielles distribuées}

Les \textbf{embeddings vectoriels} constituent une avancée fondamentale en traitement du langage naturel. Mikolov \textit{et al.}~\cite{mikolov2013} ont démontré qu'il est possible de représenter des mots, des phrases ou des documents sous forme de vecteurs denses dans un espace de dimension réduite ($\mathbb{R}^d$, typiquement $d \in [128, 1536]$), de sorte que la proximité géométrique entre vecteurs reflète la similarité sémantique entre les entités représentées.

La \textbf{similarité cosinus} est la mesure privilégiée pour comparer deux embeddings $\mathbf{a}$ et $\mathbf{b}$~:

\begin{equation}
\text{sim}(\mathbf{a}, \mathbf{b}) = \frac{\mathbf{a} \cdot \mathbf{b}}{\|\mathbf{a}\| \cdot \|\mathbf{b}\|} = \frac{\sum_{i=1}^{d} a_i b_i}{\sqrt{\sum_{i=1}^{d} a_i^2} \cdot \sqrt{\sum_{i=1}^{d} b_i^2}}
\end{equation}

\noindent Une valeur de 1 indique une identité sémantique parfaite, tandis qu'une valeur de 0 signale l'absence de relation sémantique. Les modèles de type \textbf{Sentence-BERT} ou \textbf{text-embedding} des fournisseurs LLM étendent ce principe à des séquences entières de texte.

\subsection{Recherche par plus proches voisins}

La recherche de similarité dans des espaces vectoriels de grande dimension pose un défi algorithmique majeur. Johnson, Douze et Jégou~\cite{johnson2019} ont proposé des approches de recherche approximative de plus proches voisins (\textit{Approximate Nearest Neighbor}, ANN) accélérées par GPU, permettant de parcourir des milliards de vecteurs en temps quasi-constant. Les structures d'index les plus courantes sont~:

\begin{itemize}[noitemsep]
    \item \textbf{IVF} (\textit{Inverted File Index})~: partitionnement de l'espace en cellules de Voronoï, permettant de limiter la recherche aux cellules voisines du vecteur requête~;
    \item \textbf{HNSW} (\textit{Hierarchical Navigable Small World})~: graphe de navigation hiérarchique offrant un excellent compromis précision/vitesse pour les requêtes en ligne~;
    \item \textbf{Product Quantization}~: compression des vecteurs en sous-quantificateurs, réduisant l'empreinte mémoire tout en préservant la précision de recherche.
\end{itemize}

\subsection{Bases de données vectorielles}

L'essor des LLM a engendré une nouvelle catégorie de systèmes de stockage~: les \textbf{bases de données vectorielles}. Notre plateforme utilise \textbf{pgvector}, une extension PostgreSQL qui ajoute un type natif \texttt{VECTOR(n)} et des opérateurs de recherche par similarité directement dans le SGBDR existant. Cette approche offre l'avantage de ne pas introduire un composant d'infrastructure supplémentaire, les embeddings cohabitant avec les données relationnelles dans une même base.

\subsection{Le cache sémantique}

Le \textbf{cache sémantique} représente une évolution du cache exact traditionnel. Là où un cache clé-valeur classique ne peut retourner un résultat que pour une requête strictement identique à une requête déjà mise en cache, un cache sémantique exploite les embeddings vectoriels pour identifier les requêtes \textit{sémantiquement proches} d'une requête déjà traitée. Si la similarité cosinus entre le vecteur de la requête entrante et celui d'une entrée en cache dépasse un seuil prédéfini (typiquement $\theta \geq 0{,}95$), la réponse mise en cache est retournée directement, économisant à la fois la latence et le coût de l'appel LLM. Dans notre plateforme, ce mécanisme est intégré au Gateway et permet de réduire significativement les coûts pour les prompts récurrents au sein d'un même tenant.


% ────────────────────────────────────────────────────────────
\section{Architectures multi-tenant SaaS}
\label{sec:multitenant}
% ────────────────────────────────────────────────────────────

\subsection{Définition du multi-tenant}

Une architecture \textbf{multi-tenant} (\textit{multi-locataire}) est un modèle de conception logicielle dans lequel une seule instance d'une application dessert simultanément plusieurs organisations clientes, appelées \textbf{tenants}, tout en garantissant l'isolation stricte de leurs données respectives~\cite{bezemer2010}. Ce modèle est au c\oe ur des solutions SaaS (\textit{Software as a Service}) modernes.

Bezemer et Zaidman~\cite{bezemer2010} identifient trois niveaux d'isolation dans une architecture multi-tenant~:

\begin{enumerate}[noitemsep]
    \item \textbf{Base de données partagée, schéma partagé}~: tous les tenants partagent les mêmes tables, la séparation étant assurée par une colonne \texttt{tenant\_id}~;
    \item \textbf{Base de données partagée, schéma séparé}~: chaque tenant dispose de son propre schéma au sein de la même instance de base de données~;
    \item \textbf{Base de données séparée}~: chaque tenant dispose de sa propre instance de base de données.
\end{enumerate}

Dans notre plateforme, nous avons adopté le premier niveau --- base de données partagée avec colonne \texttt{tenant\_id} --- pour sa simplicité de gestion et sa scalabilité, tout en implémentant une isolation stricte à trois niveaux (API, service, base de données).

\subsection{Contrôle d'accès basé sur les rôles (RBAC)}

Le modèle RBAC (\textit{Role-Based Access Control}) est un mécanisme de contrôle d'accès qui attribue des permissions non pas directement aux utilisateurs, mais à des rôles que les utilisateurs endossent. Conformément aux recommandations de l'OWASP~\cite{owasp2021}, notre plateforme implémente un RBAC à quatre niveaux hiérarchiques~: Super Administrateur, Administrateur Tenant, Observateur et Agent API.


% ────────────────────────────────────────────────────────────
\section{Sécurité des API~: JWT, hachage et bonnes pratiques}
\label{sec:securite}
% ────────────────────────────────────────────────────────────

La sécurisation des interfaces programmatiques constitue un enjeu critique pour toute plateforme SaaS exposée sur Internet. Cette section présente les fondements théoriques des mécanismes d'authentification et de protection déployés dans notre système.

\subsection{JSON Web Tokens (JWT)}

Le standard \textbf{JWT} (\textit{JSON Web Token}), formalisé dans la RFC~7519~\cite{rfc7519}, définit un format compact et auto-contenu pour transmettre des assertions de sécurité entre deux parties. Un JWT se compose de trois segments encodés en Base64URL et séparés par des points~:

\begin{equation}
\texttt{JWT} = \texttt{Base64(Header)}.\texttt{Base64(Payload)}.\texttt{Base64(Signature)}
\end{equation}

\begin{itemize}[noitemsep]
    \item le \textbf{Header} spécifie l'algorithme de signature (HS256, RS256) et le type de token~;
    \item le \textbf{Payload} contient les \textit{claims} --- des assertions telles que l'identité de l'utilisateur (\texttt{sub}), la date d'expiration (\texttt{exp}), le rôle (\texttt{role}) et l'identifiant du tenant (\texttt{tenant\_id})~;
    \item la \textbf{Signature} garantit l'intégrité du token par un calcul HMAC ou RSA appliqué aux deux premiers segments.
\end{itemize}

L'architecture \textbf{access/refresh token} est un pattern complémentaire recommandé par le cadre OAuth~2.0~\cite{oauth2012}. L'access token, de courte durée de vie (typiquement 15 à 30 minutes), autorise l'accès aux ressources protégées, tandis que le refresh token, de plus longue durée, permet d'obtenir un nouvel access token sans réauthentification. Cette séparation limite la fenêtre d'exposition en cas de compromission d'un access token.

\subsection{Hachage sécurisé des mots de passe}

Le stockage des mots de passe en clair constitue une vulnérabilité critique classée dans le Top~10 OWASP~\cite{owasp2021}. Les algorithmes de \textbf{hachage adaptatif} répondent à cette menace en produisant une empreinte irréversible dont le coût de calcul est volontairement élevé, rendant les attaques par force brute économiquement prohibitives.

L'algorithme \textbf{PBKDF2} (\textit{Password-Based Key Derivation Function 2}), spécifié dans la RFC~2898~\cite{rfc2898}, applique une fonction pseudo-aléatoire (typiquement HMAC-SHA256) de manière itérative~:

\begin{equation}
\texttt{DK} = \text{PBKDF2}(\texttt{PRF}, \texttt{Password}, \texttt{Salt}, c, \texttt{dkLen})
\end{equation}

\noindent où \texttt{PRF} est la fonction pseudo-aléatoire, \texttt{Salt} un sel aléatoire unique par utilisateur, $c$ le nombre d'itérations (recommandation OWASP~: $c \geq 600\,000$ pour SHA-256) et \texttt{dkLen} la longueur de la clé dérivée. Le sel empêche les attaques par tables arc-en-ciel (\textit{rainbow tables}), et le nombre élevé d'itérations impose un coût prohibitif à l'attaquant.

\subsection{Blacklisting de tokens et sécurité de session}

La révocation de tokens JWT pose un défi inhérent à leur nature stateless~: un token émis reste valide jusqu'à son expiration, même si l'utilisateur se déconnecte. La solution adoptée dans notre plateforme consiste à maintenir une \textbf{liste noire} (\textit{blacklist}) des tokens révoqués dans Redis, avec une clé TTL (\textit{Time To Live}) correspondant à la durée de vie résiduelle du token. Cette approche préserve la nature stateless de l'authentification tout en offrant un mécanisme de révocation immédiat, conformément aux recommandations OWASP pour les architectures API~\cite{owasp2021}.


% ────────────────────────────────────────────────────────────
\section{Détection d'anomalies par apprentissage automatique}
\label{sec:anomaly}
% ────────────────────────────────────────────────────────────

La détection d'anomalies constitue un domaine fondamental de l'apprentissage automatique. Chandola \textit{et al.}~\cite{chandola2009} définissent une anomalie comme une observation qui s'écarte significativement du comportement attendu dans un jeu de données. Dans le contexte de notre plateforme, les anomalies correspondent à des pics ou des creux inhabituels de consommation de tokens.

Notre système implémente \textbf{trois algorithmes complémentaires} pour maximiser la couverture de détection~:

\subsection{Décomposition STL}

La décomposition STL (\textit{Seasonal and Trend decomposition using Loess}), proposée par Cleveland \textit{et al.}~\cite{cleveland1990}, est une technique de décomposition des séries temporelles en trois composantes~:

\begin{equation}
Y_t = T_t + S_t + R_t
\end{equation}

\noindent où $Y_t$ est la valeur observée au temps $t$, $T_t$ la composante de tendance, $S_t$ la composante saisonnière et $R_t$ le résidu. Une anomalie est détectée lorsque le résidu $R_t$ dépasse un seuil défini en termes de z-scores~:

\begin{equation}
z_t = \frac{R_t - \mu_R}{\sigma_R}
\end{equation}

\noindent Un z-score $|z_t| > 3$ indique une anomalie statistiquement significative.

\subsection{Isolation Forest}

L'\textbf{Isolation Forest}, proposé par Liu \textit{et al.}~\cite{liu2008}, est un algorithme de détection d'anomalies non supervisé particulièrement efficace pour les données à haute dimensionnalité. Son principe fondamental repose sur l'observation que les anomalies, étant rares et différentes, sont plus facilement \textbf{isolables} que les observations normales.

L'algorithme construit un ensemble d'\textbf{arbres d'isolation} (\textit{isolation trees}) en sélectionnant aléatoirement un attribut et une valeur de coupure. La profondeur moyenne à laquelle une observation est isolée dans l'ensemble des arbres constitue son \textbf{score d'anomalie}~: les observations isolées rapidement (à faible profondeur) sont considérées comme anomales~\cite{liu2008}.

Le score d'anomalie est normalisé selon la formule~:

\begin{equation}
s(x, n) = 2^{-\frac{E[h(x)]}{c(n)}}
\end{equation}

\noindent où $E[h(x)]$ est la profondeur moyenne d'isolation de l'observation $x$, et $c(n)$ est la profondeur moyenne dans un arbre binaire de recherche de $n$ éléments, servant de facteur de normalisation.

\begin{figure}[H]
    \centering
    \fbox{\parbox{10cm}{\centering\vspace{3cm}\textit{Principe de l'Isolation Forest : les anomalies sont isolées en moins de coupures}\vspace{3cm}}}
    \caption{Principe de l'Isolation Forest~\cite{liu2008}}
    \label{fig:isolation_forest}
\end{figure}

\subsection{CUSUM (Cumulative Sum Control Chart)}

L'algorithme \textbf{CUSUM} (\textit{Cumulative Sum}), proposé par Page~\cite{page1954}, est une technique de contrôle statistique de processus conçue pour détecter des changements de moyenne dans une série temporelle. Il calcule une somme cumulative des écarts par rapport à une valeur de référence~:

\begin{equation}
S_t = \max(0, S_{t-1} + (x_t - \mu_0 - k))
\end{equation}

\noindent où $x_t$ est la valeur observée, $\mu_0$ la moyenne de référence et $k$ un paramètre de sensibilité. Une alarme est déclenchée lorsque $S_t$ dépasse un seuil prédéfini $h$.

L'intérêt du CUSUM dans notre contexte réside dans sa sensibilité aux \textbf{dérives progressives} de consommation, complémentaire de l'Isolation Forest qui excelle dans la détection de points isolés.

\subsection{Paradigme d'apprentissage et pipeline de données}

Une question fondamentale en apprentissage automatique concerne la nature de la supervision~: les algorithmes de notre plateforme relèvent-ils de l'apprentissage supervisé (nécessitant des données étiquetées <<~anomalie~>> / <<~normal~>>) ou de l'apprentissage non supervisé~? Dans le contexte de la gouvernance LLM, l'obtention de telles étiquettes est impraticable~: elle exigerait qu'un expert humain qualifie rétrospectivement chaque journée de consommation comme normale ou anormale pour chaque tenant. C'est pourquoi notre système repose intégralement sur un \textbf{apprentissage non supervisé} et \textbf{statistique}~\cite{chandola2009}.

\subsubsection{Phase 1~: Apprentissage hors ligne --- le Baseline}

Le Baseline constitue le <<~profil de référence~>> d'un tenant, calculé périodiquement (chaque nuit) à partir de l'historique de consommation. Le processus d'apprentissage s'effectue en quatre étapes~:

\begin{enumerate}[noitemsep]
    \item \textbf{Extraction des données.} Les agrégats quotidiens sont récupérés depuis la table \texttt{llm\_cost\_daily} sur une fenêtre glissante de 30 jours (minimum~: 14 jours). Chaque jour produit un \textbf{vecteur de caractéristiques à 6 dimensions}~:
    \begin{equation}
    \mathbf{x}_t = \begin{bmatrix} \text{total\_tokens} \\ \text{call\_count} \\ \text{avg\_tokens\_per\_call} \\ \text{cost\_usd} \\ \text{error\_rate} \\ \text{avg\_latency\_ms} \end{bmatrix}
    \end{equation}
    Ce vecteur multi-dimensionnel est essentiel, car il permet de distinguer un pic de tokens causé par davantage d'appels d'un pic causé par des prompts individuellement plus longs --- deux situations aux implications opérationnelles radicalement différentes.

    \item \textbf{Calcul des statistiques descriptives.} La moyenne $\mu$, l'écart-type $\sigma$, le minimum et le maximum des tokens quotidiens sont calculés sur la fenêtre. Ces statistiques définissent quantitativement ce qui constitue un <<~comportement normal~>> pour ce tenant spécifique.

    \item \textbf{Apprentissage structurel implicite.} L'Isolation Forest est entraîné (\texttt{fit}) sur la matrice $\mathbf{X} \in \mathbb{R}^{30 \times 6}$ des vecteurs historiques. L'algorithme construit 100 arbres d'isolation qui <<~apprennent~>> la structure de distribution des données sans jamais recevoir d'étiquettes~: il identifie les régions denses (normalité) et les régions creuses (anomalies potentielles).

    \item \textbf{Calibration adaptative du seuil.} Plutôt qu'un seuil fixe, le système calibre un seuil propre à chaque tenant en calculant le 5\textsuperscript{e} percentile des scores de l'Isolation Forest sur les données d'entraînement elles-mêmes. Ainsi, un tenant à consommation très variable aura un seuil plus permissif qu'un tenant à consommation stable.
\end{enumerate}

Le Baseline résultant est persisté dans la table \texttt{llm\_token\_baseline} et constitue la <<~mémoire apprise~>> du système.

\subsubsection{Phase 2~: Inférence en ligne --- la Détection}

À chaque nouvel appel LLM, le système compare la consommation courante au Baseline appris~:

\begin{enumerate}[noitemsep]
    \item \textbf{STL}~: décompose les 30 derniers jours + la valeur du jour en tendance + saisonnalité + résidu. Il <<~sait~>> qu'un lundi a un profil différent d'un samedi. Seul le résidu est évalué par un z-score~; un $|z| \geq 2{,}5$ déclenche un vote d'anomalie.
    \item \textbf{Isolation Forest}~: le vecteur $\mathbf{x}_{\text{today}}$ est scoré par les arbres déjà entraînés. Un score inférieur au seuil calibré signifie que le point est <<~facile à isoler~>>, donc anormal.
    \item \textbf{CUSUM}~: met à jour deux accumulateurs persistants $C^+$ et $C^-$ avec la valeur normalisée du jour. L'état des accumulateurs est sauvegardé dans le Baseline après chaque évaluation, ce qui confère à CUSUM une \textbf{mémoire inter-jours} que les deux autres algorithmes ne possèdent pas.
\end{enumerate}

\subsubsection{Vote par ensemble majoritaire}

Les trois détecteurs émettent chacun un \textbf{vote binaire} (anomalie / normal). La décision finale suit une règle de vote majoritaire~: au moins 2 votes sur 3 sont nécessaires pour confirmer une anomalie. La sévérité est ensuite graduée selon le nombre de votes et la magnitude des scores individuels~:

\begin{table}[H]
\centering
\caption{Règle de décision de l'ensemble de détection}
\label{tab:ensemble_voting}
\begin{tabularx}{\textwidth}{|c|c|X|}
\hline
\textbf{Votes concordants} & \textbf{Sévérité} & \textbf{Interprétation} \\
\hline
0--1 / 3 & Normal & Signal faible ou bruit statistique \\
\hline
2 / 3 & Warning ou High & Anomalie probable, deux algorithmes indépendants concordent \\
\hline
3 / 3 & High ou Critical & Anomalie confirmée avec très haute confiance \\
\hline
\end{tabularx}
\end{table}

Cette architecture en deux phases --- apprentissage hors ligne du Baseline, puis inférence en ligne avec vote majoritaire --- permet au système de s'auto-adapter à chaque client sans intervention humaine et sans données étiquetées, le modèle ré-apprenant continuellement à mesure que de nouvelles données sont intégrées dans la fenêtre glissante.

\subsection{Complémentarité des trois algorithmes}

La combinaison de ces trois algorithmes permet de couvrir un spectre large de types d'anomalies~:

\begin{table}[H]
\centering
\caption{Complémentarité des algorithmes de détection d'anomalies}
\label{tab:algo_complement}
\begin{tabularx}{\textwidth}{|l|X|X|}
\hline
\textbf{Algorithme} & \textbf{Forces} & \textbf{Type d'anomalies détectées} \\
\hline
STL & Capture la saisonnalité et les tendances & Déviations par rapport aux patterns temporels \\
\hline
Isolation Forest & Non supervisé, robuste aux données à haute dimension & Points aberrants isolés, clusters anormaux \\
\hline
CUSUM & Sensible aux changements graduels de moyenne & Dérives progressives, changements de régime \\
\hline
\end{tabularx}
\end{table}


% ────────────────────────────────────────────────────────────
\section{Prévision de séries temporelles}
\label{sec:forecasting}
% ────────────────────────────────────────────────────────────

Si la détection d'anomalies vise à identifier les comportements passés aberrants, la \textbf{prévision} (\textit{forecasting}) cherche à projeter les tendances futures à partir des données historiques. Dans le contexte d'une plateforme de gouvernance LLM, la prévision budgétaire permet aux décideurs d'anticiper les dépassements de coûts et de planifier les enveloppes budgétaires allouées à chaque tenant.

\subsection{Méthodes statistiques de prévision}

Makridakis, Spiliotis et Assimakopoulos~\cite{makridakis2018} ont conduit une étude comparative exhaustive des méthodes de prévision, concluant que les approches statistiques classiques surpassent souvent les modèles d'apprentissage profond pour les séries temporelles univariées de taille modérée --- précisément le profil de données rencontré dans notre contexte de consommation quotidienne de tokens.

Les méthodes les plus établies sont~:

\begin{itemize}[noitemsep]
    \item la \textbf{régression linéaire}~: modèle de base qui estime une tendance linéaire $\hat{y}_t = \alpha + \beta t$, où $\beta$ représente le taux de croissance quotidien. Sa simplicité en fait un outil de première approximation pour les projections à court terme~;
    \item les \textbf{modèles ARIMA} (\textit{AutoRegressive Integrated Moving Average})~: modèles paramétriques qui capturent à la fois les corrélations auto-régressives et les moyennes mobiles des séries, après stationnarisation par différenciation~;
    \item le \textbf{lissage exponentiel} (Holt-Winters)~: famille de modèles qui décomposent la série en niveau, tendance et saisonnalité, chaque composante étant mise à jour de manière récursive avec un poids exponentiel décroissant sur les observations passées.
\end{itemize}

\subsection{Prévision à l'échelle et analyse de scénarios}

Taylor et Letham~\cite{taylor2018} ont proposé \textbf{Prophet}, un modèle de prévision conçu pour les séries temporelles métier présentant des tendances non linéaires et des saisonnalités multiples. Le modèle décompose la série en trois composantes additives~:

\begin{equation}
y(t) = g(t) + s(t) + h(t) + \varepsilon_t
\end{equation}

\noindent où $g(t)$ est la fonction de tendance (linéaire ou logistique), $s(t)$ la composante saisonnière modélisée par une série de Fourier, $h(t)$ l'effet des jours exceptionnels, et $\varepsilon_t$ le terme d'erreur. L'intérêt de cette approche réside dans sa capacité à produire des intervalles de confiance, naturellement exploitables pour construire des scénarios multiples.

Notre plateforme adopte une approche d'\textbf{analyse de scénarios} à trois niveaux~: un scénario \textit{optimiste} (croissance freinée), un scénario \textit{moyen} (tendance linéaire prolongée) et un scénario \textit{pessimiste} (croissance accélérée). Chaque scénario est croisé avec les plafonds contractuels du tenant pour évaluer le risque de dépassement budgétaire et déclencher des alertes préventives.


% ────────────────────────────────────────────────────────────
\section{Architecture événementielle et tâches asynchrones}
\label{sec:eventdriven}
% ────────────────────────────────────────────────────────────

\subsection{Le paradigme publish/subscribe}

L'architecture \textbf{événementielle} (\textit{event-driven architecture}) est un paradigme de conception dans lequel les composants d'un système communiquent par l'émission et la consommation d'événements, plutôt que par des appels synchrones directs. Hohpe et Woolf~\cite{hohpe2003} formalisent ce paradigme à travers le pattern \textbf{Publish/Subscribe}, où un producteur publie un message sur un canal (ou \textit{topic}) et un ou plusieurs consommateurs s'y abonnent de manière découplée.

Kleppmann~\cite{kleppmann2017} souligne que ce découplage offre trois avantages fondamentaux dans les systèmes de traitement de données~:

\begin{itemize}[noitemsep]
    \item la \textbf{résilience}~: la défaillance d'un consommateur n'affecte pas le producteur ni les autres consommateurs~;
    \item la \textbf{scalabilité}~: de nouveaux consommateurs peuvent s'abonner à un flux existant sans modifier le producteur~;
    \item l'\textbf{extensibilité}~: l'ajout d'une nouvelle fonctionnalité (par exemple, la détection d'anomalies) se réduit à la souscription à un événement existant (\texttt{llm.call.completed}).
\end{itemize}

\subsection{Tâches asynchrones et orchestration}

Les opérations dont la latence est incompatible avec une réponse HTTP synchrone --- telles que l'agrégation de coûts sur des millions d'enregistrements ou l'appel à un LLM pour expliquer une anomalie --- doivent être déléguées à des \textbf{workers asynchrones}. Le pattern \textbf{Task Queue}~\cite{hohpe2003} repose sur trois composants~:

\begin{enumerate}[noitemsep]
    \item un \textbf{producteur} qui enfile une tâche dans une file de messages~;
    \item un \textbf{broker} (Redis ou RabbitMQ) qui stocke la tâche en mémoire~;
    \item un \textbf{worker} qui consomme la tâche, l'exécute et persiste le résultat.
\end{enumerate}

\textbf{Celery} est le framework de référence en Python pour l'orchestration de tâches asynchrones distribuées~\cite{celery2024}. Son module \textbf{Celery Beat} ajoute un planificateur (\textit{scheduler}) qui déclenche automatiquement des tâches périodiques selon une expression crontab. Dans notre plateforme, Celery Beat orchestre trois tâches critiques~: l'agrégation quotidienne des coûts (chaque nuit), le rollup mensuel (chaque premier du mois) et le scan d'optimisation hebdomadaire.

\subsection{Pipelines de données}

Les pipelines de données de notre plateforme suivent le paradigme \textbf{ELT} (\textit{Extract, Load, Transform})~\cite{kleppmann2017}, où les données brutes sont d'abord chargées dans la base de données (table \texttt{llm\_token\_log}), puis transformées par des tâches Celery en agrégats exploitables (\texttt{llm\_cost\_daily}, \texttt{llm\_cost\_monthly}). Cette approche est préférable au paradigme ETL traditionnel dans un contexte cloud, car elle exploite la puissance de calcul du SGBDR pour les transformations et évite les erreurs de conversion en amont.


% ────────────────────────────────────────────────────────────
\section{Conteneurisation et intégration continue}
\label{sec:conteneurisation}
% ────────────────────────────────────────────────────────────

\subsection{Principes de la conteneurisation}

La \textbf{conteneurisation} est une technique de virtualisation au niveau du système d'exploitation qui encapsule une application et ses dépendances dans une unité d'exécution isolée appelée \textbf{conteneur}~\cite{merkel2014}. Contrairement à la virtualisation traditionnelle par hyperviseur, les conteneurs partagent le noyau de l'hôte et ne transportent que les bibliothèques strictement nécessaires, offrant un démarrage quasi instantané et une empreinte mémoire minimale.

Merkel~\cite{merkel2014} identifie trois propriétés fondamentales qui justifient l'adoption massive de Docker dans l'industrie~:

\begin{itemize}[noitemsep]
    \item la \textbf{reproductibilité}~: un conteneur construit à partir d'un \texttt{Dockerfile} produit un environnement d'exécution identique, indépendamment de la machine hôte~;
    \item l'\textbf{isolation}~: chaque conteneur dispose de son propre espace de noms (\textit{namespace}) pour les processus, le réseau et le système de fichiers~;
    \item la \textbf{portabilité}~: une image Docker peut être transférée et exécutée sur tout système supportant le runtime Docker, du poste de développement au serveur de production.
\end{itemize}

\subsection{Orchestration multi-conteneurs}

Les applications modernes se composent de multiples services interdépendants. \textbf{Docker Compose} est un outil d'orchestration qui permet de décrire, dans un fichier déclaratif YAML, l'ensemble des services, réseaux et volumes d'une application multi-conteneurs~\cite{merkel2014}. Pour les architectures de plus grande envergure nécessitant de l'auto-scaling et de la haute disponibilité, \textbf{Kubernetes} propose une orchestration de niveau production avec des mécanismes de \textit{self-healing}, de \textit{rolling updates} et de \textit{service discovery}~\cite{burns2018}.

Notre plateforme utilise Docker Compose pour orchestrer six services conteneurisés --- PostgreSQL (avec TimescaleDB et pgvector), Redis, l'API FastAPI, le worker Celery, le scheduler Celery Beat et Nginx --- au sein d'un réseau interne isolé avec des volumes persistants pour les données.

\subsection{Intégration continue et déploiement continu (CI/CD)}

L'\textbf{intégration continue} (\textit{Continuous Integration}, CI) est une pratique d'ingénierie logicielle qui consiste à fusionner fréquemment le code dans un tronc commun, chaque fusion déclenchant automatiquement une suite de vérifications~: analyse syntaxique (\textit{linting}), exécution des tests unitaires et d'intégration, et construction de l'image Docker. Le \textbf{déploiement continu} (\textit{Continuous Deployment}, CD) prolonge cette chaîne en automatisant la mise en production des artefacts validés~\cite{ford2017}.

Notre pipeline CI/CD, implémenté via \textbf{GitHub Actions}, automatise quatre étapes à chaque push~: la vérification syntaxique de tous les fichiers Python, l'exécution des 180 tests avec des services TimescaleDB et Redis éphémères, la construction et le push de l'image Docker multi-stage vers GitHub Container Registry (GHCR), et la notification de déploiement. Cette automatisation s'inscrit dans les principes \textbf{GitOps} décrits dans le chapitre~1, où l'état souhaité de l'infrastructure est entièrement versionné dans le dépôt Git.


% ────────────────────────────────────────────────────────────
\section{Solutions existantes et analyse comparative}
\label{sec:comparatif}
% ────────────────────────────────────────────────────────────

Plusieurs solutions de monitoring et de gouvernance des LLM existent sur le marché. Nous présentons ici une analyse comparative des principaux acteurs~:

\begin{table}[H]
\centering
\caption{Analyse comparative des solutions existantes}
\label{tab:comparatif}
\begin{tabularx}{\textwidth}{|p{2.2cm}|X|X|X|}
\hline
\textbf{Solution} & \textbf{Points forts} & \textbf{Limites} & \textbf{Différenciation LLM Dashboard} \\
\hline
LangSmith (LangChain) & Traçage détaillé des chaînes LLM, évaluation & Pas de gouvernance, pas de facturation, couplé à LangChain & Gouvernance temps réel + facturation intégrée + indépendant du framework \\
\hline
Helicone & Proxy léger, analytics de base & Pas de détection d'anomalies ML, pas de multi-tenant complet & 3 algorithmes ML + isolation multi-tenant + RBAC 4 niveaux \\
\hline
PortKey AI & Gateway multi-fournisseur, cache sémantique & Pas de facturation contractuelle, pas de prévision budgétaire & Contrats \textit{pay-as-you-go}/forfait/hybride + prévision à 3 scénarios \\
\hline
LiteLLM & Proxy open-source, 100+ modèles & Pas d'UI, pas d'analytics, pas de détection d'anomalies & Dashboard exécutif complet + WebSocket temps réel + assistant IA \\
\hline
\end{tabularx}
\end{table}

Notre plateforme se distingue par son approche \textbf{intégrée de bout en bout}, combinant dans un seul système le proxying, la gouvernance, la détection d'anomalies par ML, l'explication IA, la prévision budgétaire, la facturation contractuelle et un tableau de bord exécutif interactif --- une combinaison qu'aucune solution existante ne propose de manière unifiée.


\section*{Conclusion}

Ce chapitre a établi un socle théorique complet pour les choix architecturaux et algorithmiques de notre plateforme. L'étude de l'architecture Transformer et des mécanismes de tokenisation a fondé la compréhension du modèle économique des LLM et la nécessité d'un monitoring granulaire. L'analyse des enjeux de gouvernance et du pattern API Gateway a justifié le positionnement d'un proxy intelligent entre les applications clientes et les fournisseurs de modèles. La présentation des séries temporelles, de TimescaleDB et des hypertables a posé les bases de notre système d'agrégation de coûts, tandis que l'étude des embeddings vectoriels et du cache sémantique a introduit les mécanismes d'optimisation exploités par le module Gateway et le module Explainer. Les sections consacrées à la sécurité API (JWT, PBKDF2), à l'architecture multi-tenant et au RBAC ont défini le cadre de protection des données et d'isolation des tenants. Les algorithmes de détection d'anomalies (STL, Isolation Forest, CUSUM) et les méthodes de prévision budgétaire fournissent les fondements mathématiques des modules d'intelligence artificielle. Enfin, l'étude de l'architecture événementielle, des tâches asynchrones Celery, de la conteneurisation Docker et de l'intégration continue CI/CD a complété le panorama des paradigmes d'ingénierie logicielle mis en \oe uvre. L'analyse comparative des solutions existantes a confirmé l'originalité de notre approche intégrée de bout en bout. Le chapitre suivant détaille le Sprint~0, consacré à l'identification des besoins et à la définition de l'environnement de travail.
